% Rozwinięcie TO-18 i AM II-01, przygotowanie do szeregów funkcyjnych.
\chapter{Zbieżność punktowa i jednostajna}
\label{ma2:ch:uniform-convergence}

Niech $X$ będzie zbiorem, a $(Y,d)$ przestrzenią
metryczną. Gdy $X$ jest przestrzenią topologiczną,
przez $C(X,Y)$ oznaczamy zbiór wszystkich ciągłych
odwzorowań $X\to Y$. Rozróżnimy sposoby zbieżności ciągów
odwzorowań $X\to Y$. W przypadku przestrzeni
unormowanych wykorzystamy wyniki rozdziału
\ref{ma2:ch:normed}. Topologiczne ujęcie przestrzeni
odwzorowań znajduje się w topologii ogólnej, rozdział 18.

\section{Definicje i pierwsze kryteria}

\begin{definicja}
\label{ma2:def:pointwise-uniform}
Ciąg $(f_n)_{n\ge1}$ odwzorowań $X\to Y$ jest
zbieżny punktowo do $f:X\to Y$, gdy
\[
 \forall x\in X\ \forall\varepsilon>0\ \exists N\ge1\
 \forall n\ge N:\quad d(f_n(x),f(x))<\varepsilon.
\]
Jest zbieżny jednostajnie do $f$ na $X$, gdy
\[
 \forall\varepsilon>0\ \exists N\ge1\
 \forall n\ge N\ \forall x\in X:\quad
 d(f_n(x),f(x))<\varepsilon.
\]
\end{definicja}
Różnica polega na tym, że w drugim przypadku wybór
$N$ jest niezależny od $x$. Zbieżność jednostajna
implikuje zbieżność punktową. Granica jest jednoznaczna
z jednoznaczności granicy ciągu w przestrzeni metrycznej
(AM I, rozdział 06, twierdzenie o jednoznaczności
granicy, etykieta \texttt{ma1:thm:unique-limit}).
Istotnie, dwie granice punktowe muszą mieć tę samą
wartość w każdym punkcie dziedziny. Na zbiorze pustym obie własności
zachodzą automatycznie; istnieje dokładnie jedno
odwzorowanie z tego zbioru do $Y$.

\begin{twierdzenie}[Kryteria zbieżności jednostajnej]
\label{ma2:thm:uniform-criteria}
Niech $X\ne\varnothing$, a $f_n,f:X\to Y$
będą odwzorowaniami. Następujące warunki są równoważne:
\begin{enumerate}[label=(\roman*)]
\item $f_n\to f$ jednostajnie na $X$;
\item $M_n:=\sup_{x\in X}d(f_n(x),f(x))\to0$,
gdzie dopuszczamy $M_n=+\infty$;
\item dla każdego ciągu $(x_n)$ w $X$ mamy
$d(f_n(x_n),f(x_n))\to0$.
\end{enumerate}
\end{twierdzenie}
\begin{proof}
Z (i), zastosowanego z dokładnością $\varepsilon/2$,
mamy $M_n\le\varepsilon/2<\varepsilon$ od pewnego
indeksu, co daje (ii). Z (ii) wynika (i), bo
$d(f_n(x),f(x))\le M_n$ dla każdego $x$.
Z (ii) otrzymujemy również (iii): dla dowolnego
ciągu $(x_n)$ od pewnego indeksu mamy
$0\le d(f_n(x_n),f(x_n))\le M_n$, a $M_n\to0$.
Gdyby (i) nie zachodziło, istniałoby $\varepsilon>0$
takie, że dla dowolnego $N$ można wybrać $n\ge N$
i $x\in X$ z $d(f_n(x),f(x))\ge\varepsilon$.
Wybieramy rekurencyjnie ściśle rosnące indeksy
$n_k$ i punkty $x_{n_k}$ spełniające tę nierówność:
najpierw stosujemy warunek z $N=1$, a po wyborze
$n_k$ stosujemy go z $N=n_k+1$.
Uzupełniamy pozostałe wyrazy ciągu $(x_n)$ dowolnym
ustalonym punktem $X$. Wzdłuż podciągu $(n_k)$
odległości są nie mniejsze niż $\varepsilon$,
co przeczy (iii).
\end{proof}

\begin{przyklad}
\label{ma2:ex:powers-pointwise}
Na $[0,1]$ funkcje $f_n(x)=x^n$ zbiegają punktowo
do funkcji $f$ równej $0$ na $[0,1)$ i $1$ w punkcie $1$.
Dla $x<1$ wynika to ze zbieżności potęg liczby
o module mniejszym od $1$ (AM I, rozdział 06, przykład ciągu geometrycznego,
etykieta \texttt{ma1:ex:geometric-sequence}).
Dla każdego $n$ mamy $\sup_{[0,1]}|f_n-f|=1$:
na $[0,1)$ wartości $x^n$ są mniejsze od $1$,
ale dążą do $1$ przy $x\to1^-$, a w punkcie $1$
różnica wynosi $0$. Zatem zbieżność nie jest
jednostajna. Na $[0,r]$, gdzie $0\le r<1$,
supremum różnicy wynosi $r^n\to0$, więc tam
zbieżność jest jednostajna. Rysunek
\ref{ma2:fig:powers-uniform} ilustruje oba przypadki.
\end{przyklad}

\begin{figure}[htbp]
\centering
\begin{tikzpicture}[x=8cm,y=3.8cm,>=Stealth]
 \draw[->] (-.035,0)--(1.08,0) node[right]{$x$};
 \draw[->] (0,-.06)--(0,1.14) node[above]{$y$};
 \draw[gray!55,densely dotted] (0,1)--(1,1)--(1,0);
 \foreach \n/\kol in {1/blue!75!black,2/teal!80!black,4/orange!90!black,12/red!70!black}{
  \draw[\kol,thick,domain=0:1,samples=81] plot (\x,{pow(\x,\n)});
 }
 \draw[black,very thick] (0,0)--(.992,0);
 \draw[black,fill=white] (1,0) circle[radius=1.8pt];
 \fill (1,1) circle[radius=1.8pt];
 \draw[gray,dashed] (.7,0)--(.7,.7);
 \node[below] at (.7,0) {$r$};
 \node[below] at (1,0) {$1$};
 \node[left] at (0,1) {$1$};
 \node[below left] at (0,0) {$0$};
 \node[blue!75!black,anchor=west] at (.05,.87) {$n=1$};
 \node[teal!80!black,anchor=west] at (.05,.73) {$n=2$};
 \node[orange!90!black,anchor=west] at (.05,.59) {$n=4$};
 \node[red!70!black,anchor=west] at (.05,.45) {$n=12$};
\end{tikzpicture}
\caption{Funkcje $x^n$ i ich granica punktowa (czarna).
Na każdym $[0,r]$, $r<1$, błąd nie przekracza $r^n$,
lecz na całym $[0,1]$ jego supremum wynosi $1$.}
\label{ma2:fig:powers-uniform}
\end{figure}
\FloatBarrier

\begin{przyklad}[Ciągła granica punktowa nie wystarcza]
\label{ma2:ex:moving-peak}
Na $[0,1]$ połóżmy
$g_n(x)=\max\{1-n|x-1/n|,0\}$.
Funkcje są ciągłe, $g_n(0)=0$, a dla ustalonego
$x>0$ mamy $g_n(x)=0$ od chwili, gdy $2/n<x$.
Zatem granicą punktową jest funkcja stale równa $0$.
Jednak $g_n(1/n)=1$ i $0\le g_n\le1$, więc
$\sup_{[0,1]}|g_n|=1$. Zbieżność nie jest jednostajna,
mimo ciągłości granicy i zwartości dziedziny.
\end{przyklad}

\section{Jednostajne kryterium Cauchy'ego}

\begin{twierdzenie}
\label{ma2:thm:uniform-cauchy}
Niech $(f_n)$ będzie ciągiem odwzorowań $X\to Y$.
Jeśli $f_n\to f$ jednostajnie, to
\[
 \forall\varepsilon>0\ \exists N\ge1\
 \forall m,n\ge N\ \forall x\in X:\quad
 d(f_m(x),f_n(x))<\varepsilon.
\]
Jeżeli $Y$ jest zupełna, warunek ten jest także
wystarczający do zbieżności jednostajnej do pewnej
funkcji $f:X\to Y$.
\end{twierdzenie}
\begin{proof}
Dla $\varepsilon>0$ wybieramy $N$ z definicji
zbieżności jednostajnej dla dokładności $\varepsilon/2$.
Dla $m,n\ge N$ i wszystkich $x\in X$ mamy
\[
 d(f_m(x),f_n(x))
 \le d(f_m(x),f(x))+d(f(x),f_n(x))<\varepsilon.
\]
To daje warunek konieczny.
Odwrotnie, dla każdego $x$ ciąg $(f_n(x))$ jest
ciągiem Cauchy'ego, więc z zupełności $Y$ ma
jednoznaczną granicę $f(x)$.
Dla $\varepsilon>0$ wybieramy wspólny $N$
w warunku Cauchy'ego z dokładnością $\varepsilon/2$.
Dla $n\ge N$ i każdego $x$, przechodząc przy ustalonych
$n,x$ do granicy $m\to\infty$, otrzymujemy
\[
 d(f_n(x),f(x))\le\varepsilon/2<\varepsilon.
\]
Przejście do granicy uzasadnia nierówność
$|d(y,z)-d(y,w)|\le d(z,w)$.
Indeks $N$ jest niezależny od $x$, co daje tezę.
\end{proof}
Jeżeli granica punktowa $f:X\to Y$ już istnieje,
to jednostajny warunek Cauchy'ego daje zbieżność
jednostajną do $f$ także bez zupełności $Y$:
w powyższym dowodzie zupełność służyła wyłącznie
do uzyskania granic punktowych. Wystarczy przejść
do granicy po $m$ z tym samym wspólnym progiem $N$.

Jeżeli $Y$ nie jest zupełna, wybierzmy w niej ciąg
Cauchy'ego $(y_n)$ bez granicy. Funkcje stałe
$f_n(x)=y_n$ na niepustym $X$ spełniają warunek
jednostajny Cauchy'ego, ale nie mają nawet granicy
punktowej. Zupełność jest zatem istotna.

\begin{twierdzenie}[Metryka zbieżności jednostajnej]
\label{ma2:thm:uniform-metric}
Na zbiorze $Y^X$ wszystkich odwzorowań $X\to Y$
wzór
\[
 D(f,g)=\sup_{x\in X}\min\{1,d(f(x),g(x))\}
\]
określa metrykę, przy konwencji $D=0$ dla $X=\varnothing$.
Zbieżność w tej metryce jest zbieżnością jednostajną.
Jeżeli $X\ne\varnothing$, metryka $D$ jest zupełna
wtedy i tylko wtedy, gdy $Y$ jest zupełna.
\end{twierdzenie}
\begin{proof}
Wartości $D$ należą do $[0,1]$. Symetria
i nieujemność wynikają z aksjomatów metryki $d$.
Jeśli $D(f,g)=0$, to dla każdego $x$ mamy
$\min\{1,d(f(x),g(x))\}=0$, czyli $f(x)=g(x)$.
Odwrotnie, $f=g$ daje $D(f,g)=0$.
Dla $a,b\ge0$ zachodzi
$\min\{1,a+b\}\le\min\{1,a\}+\min\{1,b\}$:
jeśli któryś składnik jest co najmniej $1$, prawa
strona jest co najmniej $1$, a w przeciwnym przypadku
wynosi $a+b$. Stosując tę nierówność do nierówności
trójkąta dla $d$ i biorąc supremum, otrzymujemy
nierówność trójkąta dla $D$.
Dla $0<\varepsilon<1$, jeśli $D(f_n,f)<\varepsilon$,
to $d(f_n(x),f(x))<\varepsilon$ dla wszystkich $x$.
Odwrotnie, jednostajna zbieżność z dokładnością
$\varepsilon/2$ daje $D(f_n,f)\le\varepsilon/2$.
Dla dowolnych $\varepsilon>0$ wystarczy w tych
argumentach użyć dokładności mniejszej od
$\min\{1,\varepsilon\}$, co obejmuje również
$\varepsilon\ge1$.
Te same oszacowania zastosowane do par $f_m,f_n$
pokazują równoważność warunku Cauchy'ego w metryce
$D$ i jednostajnego warunku Cauchy'ego.
Z twierdzenia \ref{ma2:thm:uniform-cauchy}
wynika zupełność, gdy $Y$ jest zupełna.
Odwrotnie, dla ciągu Cauchy'ego $(y_n)$ w $Y$
bierzemy funkcje stałe $f_n(x)=y_n$. Są ciągiem
Cauchy'ego w $D$, zatem przy zupełności $D$ mają
granicę $f$. Ewaluacja w dowolnym $x_0\in X$
daje $y_n\to f(x_0)$, czyli zupełność $Y$.
Dla pustego $X$ przestrzeń funkcji jest jednoelementowa.
\end{proof}

W przypadku przestrzeni unormowanej $E$ i funkcji
ograniczonych na niepustym $X$ można używać normy supremum:
\[
 \|f_n-f\|_\infty=\sup_{x\in X}\|f_n(x)-f(x)\|.
\]
Zbieżność w tej normie jest zbieżnością jednostajną
(z twierdzenia \ref{ma2:thm:uniform-criteria}).
Jednostajna granica funkcji ograniczonych jest
ograniczona: wybieramy $n$ z
$\sup_x\|f_n(x)-f(x)\|\le1$, skąd
$\|f(x)\|\le\|f_n\|_\infty+1$ dla każdego $x$.
Dla $X=\varnothing$ przestrzeń funkcji ograniczonych
jest przestrzenią zerową z normą $0$.
Zupełność $B(X,E)$ dla niepustego $X$ jest równoważna
zupełności $E$ na mocy twierdzenia \ref{ma2:thm:bounded}.
Nie powtarzamy tego dowodu.

\section{Ciągłość granicy jednostajnej}

\begin{twierdzenie}
\label{ma2:thm:uniform-continuous-limit}
Niech $X$ będzie przestrzenią topologiczną,
$(Y,d)$ przestrzenią metryczną i $f_n:X\to Y$
ciągłymi odwzorowaniami. Jeżeli $f_n\to f$
jednostajnie, to $f$ jest ciągła.
Jeśli $X$ jest metryczna i wszystkie $f_n$ są
jednostajnie ciągłe, to $f$ jest jednostajnie ciągła.
\end{twierdzenie}
\begin{proof}
Pierwsza część uogólnia argument z dowodu
twierdzenia \ref{ma2:thm:continuous} na dowolną
dziedzinę topologiczną i przeciwdziedzinę metryczną.
Ustalmy $x_0\in X$ i $\varepsilon>0$.
Wybieramy $n$ tak, aby dla każdego $x\in X$
$d(f_n(x),f(x))<\varepsilon/3$.
Ciągłość $f_n$ zapewnia otwarte otoczenie $U$
punktu $x_0$ z $d(f_n(x),f_n(x_0))<\varepsilon/3$
dla $x\in U$. Dla tych $x$ mamy
\[
 d(f(x),f(x_0))
 \le d(f(x),f_n(x))+d(f_n(x),f_n(x_0))
       +d(f_n(x_0),f(x_0))<\varepsilon.
\]
To dowodzi ciągłości.
W przypadku jednostajnej ciągłości wybieramy
$n$ tak samo, po czym bierzemy $\delta>0$
z jednostajnej ciągłości $f_n$ dla
$\varepsilon/3$. Dla dowolnych $x,y\in X$
o odległości mniejszej od $\delta$ to samo
oszacowanie, z $y$ zamiast $x_0$, daje
$d(f(x),f(y))<\varepsilon$.
Wybór $\delta$ jest niezależny od $x,y$.
\end{proof}
Przykład \ref{ma2:ex:powers-pointwise} pokazuje,
że granica punktowa funkcji ciągłych może
być nieciągła.

\begin{wniosek}
\label{ma2:cor:uniform-complete-continuous}
Jeśli $X$ jest przestrzenią topologiczną, a $Y$
zupełną przestrzenią metryczną, to $C(X,Y)$
jest zupełna w metryce $D$ z twierdzenia
\ref{ma2:thm:uniform-metric}.
\end{wniosek}
\begin{proof}
Ciąg Cauchy'ego w tej metryce ma granicę
jednostajną w $Y^X$ z twierdzenia
\ref{ma2:thm:uniform-metric}. Granica jest ciągła
z twierdzenia \ref{ma2:thm:uniform-continuous-limit}.
\end{proof}

\section{Zbieżność jednostajna na zbiorach zwartych}

\begin{definicja}
\label{ma2:def:compact-uniform}
Niech $X$ będzie przestrzenią topologiczną.
Ciąg $(f_n)$ odwzorowań $X\to Y$ jest zbieżny
jednostajnie na zbiorach zwartych do $f$,
gdy dla każdego zwartego $K\subseteq X$
ograniczenia $f_n|_K$ zbiegają jednostajnie do $f|_K$.
\end{definicja}
Zbieżność jednostajna na całym $X$ implikuje
zbieżność jednostajną na zbiorach zwartych,
a ta implikuje zbieżność punktową (zbiór
jednoelementowy jest zwarty). Jeśli $X$ jest
zwarta, pierwsze dwie własności są równoważne.
Jeśli $X$ jest skończonym zbiorem, zbieżność
punktowa implikuje zbieżność jednostajną:
dla każdego punktu wybieramy próg dla ustalonej
dokładności i bierzemy maksimum skończonej
liczby progów; przypadek pusty jest automatyczny.

\begin{przyklad}[Zbieżność na zbiorach zwartych nie jest globalnie jednostajna]
\label{ex:topog-odwzorowania-gorki-uciekajace}
Przytaczamy przykład z topologii ogólnej, rozdział 18,
sekcja 2. Na $\R$ niech
$f_n(x)=\max\{1-|x-n|,0\}$.
Każda funkcja jest ciągła, $\sup_{\R}|f_n|=1$.
Dla dowolnego zwartego $K\subseteq\R$ istnieje $M$
z $K\subseteq[-M,M]$. Gdy $n>M+1$, funkcja
$f_n$ znika na $K$. Stąd zbieżność jednostajna
na zbiorach zwartych do zera, lecz nie na całej $\R$.
\end{przyklad}
Przykład \ref{ma2:ex:moving-peak} pokazuje natomiast,
że zbieżność punktowa nie implikuje zbieżności
jednostajnej na zbiorach zwartych, nawet przy
ciągłej granicy.

\begin{twierdzenie}
\label{ma2:thm:compact-uniform-continuity}
Niech $X,Y$ będą przestrzeniami metrycznymi.
Jeżeli $(f_n)$ jest ciągiem ciągłych odwzorowań
$X\to Y$ zbieżnym do $f$ jednostajnie na zbiorach
zwartych, to $f$ jest ciągła.
\end{twierdzenie}
\begin{proof}
Niech $x_j\to x$ w $X$. Zbiór
$K=\{x\}\cup\{x_j:j\ge1\}$ jest zwarty:
w dowolnym pokryciu otwartym element zawierający
$x$ zawiera wszystkie $x_j$ od pewnego indeksu,
a pozostałe punkty można pokryć skończoną liczbą
elementów pokrycia. Na $K$ mamy zbieżność
jednostajną funkcji ciągłych, więc
twierdzenie \ref{ma2:thm:uniform-continuous-limit}
daje ciągłość $f|_K$. W szczególności
$f(x_j)\to f(x)$. Ponieważ ciąg był dowolny,
kryterium ciągowe ciągłości (AM I, rozdział 15,
twierdzenie o kryterium ciągowym, etykieta
\texttt{ma1:thm:continuity-sequential}) daje ciągłość $f$.
\end{proof}
Na przestrzeniach lokalnie zwartych Hausdorffa
można również wykorzystać zwarte otoczenia
(topologia ogólna, rozdział 14). Dla każdego $x$
wybieramy zwarty zbiór $K$ zawierający otwarte
otoczenie $U$ punktu $x$. Granica $f|_K$ jest
ciągła z twierdzenia \ref{ma2:thm:uniform-continuous-limit},
a zatem $f|_U$ jest ciągła. Ponieważ $U$ jest
otwarty w $X$, daje to ciągłość $f$ w $x$.

Przypomnijmy, że topologia zwarto-otwarta na $C(X,Y)$
ma podbazę zbiorów
\[
 [K,V]=\{g\in C(X,Y):g(K)\subseteq V\},
\]
gdzie $K\subseteq X$ jest zwarty, a $V\subseteq Y$
otwarty. Topologia zbieżności jednostajnej na
zbiorach zwartych ma w punkcie $f$ bazę otoczeń
\[
 B_K(f,\varepsilon)=
 \{g\in C(X,Y):d(f(x),g(x))<\varepsilon
                    \text{ dla wszystkich }x\in K\}.
\]
Są to definicje z topologii ogólnej, rozdział 18,
sekcje 1 i 2. Dla niepustego $K$ odległość
$d(f(x),g(x))$ osiąga maksimum na $K$, ponieważ
jest ciągła; dlatego warunek punktowy z ustalonym
$\varepsilon$ jest tu równoważny ostrej nierówności
dla supremum. Dla $K=\varnothing$ otrzymujemy całe
$C(X,Y)$.

\begin{twierdzenie}[Równoważność dwóch topologii]
\label{thm:topog-odwzorowania-zwarto-otwarta-jednostajna}
Na $C(X,Y)$ dla dowolnej przestrzeni topologicznej
$X$ i metrycznej $Y$ topologia zwarto-otwarta
jest równa topologii zbieżności jednostajnej
na zbiorach zwartych. W szczególności ciąg
funkcji ciągłych zbiega w tej topologii do
$f\in C(X,Y)$ dokładnie wtedy, gdy zbiega do
$f$ jednostajnie na zbiorach zwartych.
\end{twierdzenie}
\emph{Dowód:} topologia ogólna, rozdział 18, sekcja 2,
twierdzenie o równoważności dwóch topologii
(z tą samą etykietą źródłową) i następujący po nim
wniosek o zbieżności sieci.
Dla lokalnie zwartej, $\sigma$-zwartej przestrzeni
Hausdorffa metryzacja i zupełność tej topologii
przy zupełnym $Y$ zostały udowodnione w sekcji 4
tego samego rozdziału. Nie powtarzamy tych dowodów.

\section{Twierdzenie Diniego}

\begin{twierdzenie}[Twierdzenie Diniego]
\label{ma2:thm:dini}
Niech $X$ będzie zwartą przestrzenią topologiczną,
a $(f_n)$ ciągiem ciągłych funkcji rzeczywistych
na $X$ zbieżnym punktowo do ciągłej funkcji $f$.
Jeżeli dla każdego $x\in X$ zachodzi
$f_n(x)\le f_{n+1}(x)$ dla wszystkich $n$,
albo dla każdego $x\in X$ zachodzi
$f_n(x)\ge f_{n+1}(x)$ dla wszystkich $n$,
to zbieżność jest jednostajna.
\end{twierdzenie}
\begin{proof}
W przypadku rosnącym mamy $f_n(x)\le f(x)$
dla każdego $n,x$, ponieważ ciąg rzeczywisty
jest rosnący i zbieżny do $f(x)$.
Dla $\varepsilon>0$ definiujemy otwarte zbiory
\[
 U_n=\{x\in X:f(x)-f_n(x)<\varepsilon\}.
\]
Są otwarte dzięki ciągłości różnicy,
rosnące ze względu na inkluzję i pokrywają $X$
z punktowej zbieżności. Jeżeli $X\ne\varnothing$,
zwartość daje skończone podpokrycie
$U_{n_1},\ldots,U_{n_s}$. Biorąc $N=\max_j n_j$,
otrzymujemy $X=U_N$ dzięki rosnącemu układowi.
Dla $n\ge N$ i wszystkich $x$ zachodzi więc
$0\le f(x)-f_n(x)<\varepsilon$.
Jest to zbieżność jednostajna.
W przypadku malejącym stosujemy ten sam dowód
do $-f_n$ i $-f$; przypadek pustej dziedziny
wynika z definicji.
\end{proof}

\begin{uwaga}
W istocie argument Diniego działa także dla
odwzorowań ciągłych $f_n,f:X\to(Y,d)$, gdzie $X$
jest zwarta, jeśli $f_n\to f$ punktowo i dla każdego
$x$ ciąg $d(f_n(x),f(x))$ jest nierosnący.
Zbiory $\{x:d(f_n(x),f(x))<\varepsilon\}$ są otwarte,
gdyż nierówność
\[
 |d(y,z)-d(y',z')|\le d(y,y')+d(z,z')
\]
wynikająca z dwukrotnego zastosowania nierówności
trójkąta daje ciągłość odległości względem obu
argumentów. Zbiory są rosnące dzięki założeniu
o nierosnących odległościach.
Pokrywają $X$ z punktowej zbieżności, więc argument
skończonego podpokrycia z dowodu twierdzenia
\ref{ma2:thm:dini} daje zbieżność jednostajną.
Dla funkcji rzeczywistych obejmuje to również
przypadek, w którym kierunek monotoniczności
ciągu $f_n(x)$ zależy od $x$.
\end{uwaga}

Założenia twierdzenia są istotne. Na niezwartym
$[0,1)$ funkcje $x^n$ maleją punktowo do ciągłego
zera, ale supremum $x^n$ nadal wynosi $1$.
Na $[0,1]$ ten sam ciąg ma nieciągłą granicę
(przykład \ref{ma2:ex:powers-pointwise}).
Bez monotoniczności nawet na zwartym przedziale
ciągłość granicy nie wystarcza, jak pokazuje
przykład \ref{ma2:ex:moving-peak}.

W następnym rozdziale zastosujemy kryteria
zbieżności jednostajnej do szeregów funkcyjnych.
